\section{The Contract and Its Execution Boundary}\label{sec:r37-institution}
\subsection{A saturated renewal agreement}
At the renewal date, a public branch $j\in\{1,\ldots,k\}$ is observed with probability $\pi_j>0$, where $\sum_j\pi_j=1$. Its continuation cap is $0<b_1<\cdots<b_k<1$. The provider chooses an intermediate tier $Y_j\in[0,1]$ and a terminal tier $C_j\in[0,1]$. All branches subsequently reach the same terminal public vertex. Unit payment coefficients and unit discounting make total remaining payment $Y_j+C_j$. A fixed service-benefit continuation minus the customer's outside continuation gives the cap. At the root the provider has committed to
\begin{equation}\label{eq:r37-saturation}
 \sum_j\pi_j\E[Y_j+C_j]=\bar b,
 \qquad \bar b:=\sum_j\pi_jb_j.
\end{equation}
Under expected participation, $\E[Y_j+C_j]\leq b_j$ for every $j$, with expectation conditional on that observed branch and taken before the internal draw. Positivity of the probabilities and \eqref{eq:r37-saturation} imply equality separately on every branch. Under pathwise participation, the stronger constraint is $Y_j+C_j\leq b_j$ almost surely. The same argument then implies equality almost surely, not merely in expectation.

The operating payoff is $f(C_j)-h_j(Y_j)$, with increasing concave $f$ and convex nondecreasing $h_j$ satisfying $h_j(0)=0$. General shape assumptions suffice for the deterministic and pathwise results. The global continuous algorithm specializes to the quadratic primitives in \eqref{eq:r34-primitives}. The negative intermediate contribution represents the net cost of an advance or expedited tier, while the terminal tier contributes net operating return. It does not mean that the customer has negative service benefit: that benefit has already entered the continuation cap. There are no endogenous public transitions, coupled commitments, or switching costs in the renewal theorem.

\subsection{What a writable code means}
At the observed branch, an encoder selects a symbol $Z\in\{1,\ldots,m\}$. At the terminal vertex the decoder observes only the current public vertex and $Z$. Its stored program can contain arbitrary read-only model constants but cannot recover the realized branch from those constants. An inherited intermediate tier, a customer identifier with a branch lookup, a retained shared random seed, and a transmitted numerical payment are all additional history channels; none is free in this interface. An $m$-symbol alphabet requires $\lceil\log_2m\rceil$ fixed-width writable bits, not $m$ bits. Public-state storage and read-only program storage are counted separately.

Private encoder randomization is allowed under expected participation. Fresh decoder randomness is independent of earlier private choices conditional on the symbol. Any seed that correlates those choices across dates must itself be encoded in the charged symbol. Concavity permits a deterministic terminal decoder at the expected-participation optimum, so a codebook is a collection of terminal levels. The lower bound for exact implementation continues to allow randomized encoders and decoders under this accounting. Distinct full-information terminal actions and strict concavity remain visible assumptions; the exact threshold is not claimed for repeated caps or nonunique full optima.

\subsection{An operational mechanism, not an empirical calibration}
Consider a stipulated service-batch agreement administered by a renewal desk and executed by a shared dispatch gateway. The desk observes the customer's current outside alternative and approved continuation. It can perform an advance tier and issue one of a small number of dispatch tickets. The gateway implements the terminal tier named by that ticket; its interface does not retain the renewal record. A separate billing audit can retain the contract without making that record an input to the gateway. Limiting tickets can price gateway-state certification, transmission, or a restricted instruction set. This is a documented stylized mechanism, not a claim that a particular industry already uses this architecture.

In the expected-participation version, the offer communicates the advance tier, the terminal ticket levels, and their probabilities. The customer evaluates expected remaining payment before the ticket is drawn and either renews or exits at that review. The protocol commits to the advertised draw. After renewal the desk performs the advance tier, draws the ticket, and sends only that ticket to the gateway. The realized ticket may be disclosed to the customer. ``Private'' means that the seed is neither a pre-acceptance conditioning variable nor a free future decoder state; it need not mean that realized service is hidden forever. The customer does not receive a new costless exit option after observing an unfavorable draw within this modeled renewal episode.

This version is meaningful only for an ex ante willingness-to-pay or expected-obligation cap and an agreement that permits the stated realization variation. It is not a way to satisfy an ex post enforceable payment ceiling in expectation. A hard billing limit, a service-safety requirement on every draw, or a review after disclosure instead calls for pathwise feasibility, with the deterministic frontier as its optimum. We make no legal assertion about enforceability. The two feasible sets are explicitly different design alternatives.

Repeated customers who may condition a later renewal on past tickets require that history to enter the later public state, cap, or charged memory. The three-date theorem is not an unqualified repeated-game result. Similarly, a contractible seed disclosed before acceptance changes the conditioning information and hence the participation test. It cannot simply be added as a free source of branch-dependent execution. Section~\ref{sec:r37-resources} gives a bounded-realization variant for an operator who permits some variation but not the full expected-participation lottery.
